\documentclass[parskip=full,11pt]{scrartcl}
\usepackage[utf8]{inputenc}
\usepackage[T1]{fontenc}
\usepackage[german]{babel}
\usepackage[useregional]{datetime2}
\usepackage[pdfborderstyle={/S/U/W 0}]{hyperref}
\usepackage[nameinlink]{cleveref}
\usepackage[section]{placeins}
\usepackage[top=2.5cm, bottom=2.5cm, left=4cm, right=3cm]{geometry}
\usepackage{xcolor}
\usepackage{graphicx}
\usepackage{csquotes}
\usepackage{amsmath} % for $\text{}$
\usepackage{changelog} % local .sty file
\usepackage{enumitem}
\usepackage{algorithm}
\usepackage{algorithmicx}
\usepackage{algpseudocode}

\setlist{nosep}

\newcommand\urlpart[2]{$\underbrace{\text{\texttt{#1}}}{\text{#2}}$}
\raggedbottom
\crefname{figure}{Abb}{Abb}

\newcommand\producttitle{treff.}
\newcommand\protocolversion{0.3}
\hypersetup{
	pdftitle={\producttitle~- Änderungsdokument},
	bookmarks=true,
}

% section numbers in margins:
\renewcommand\sectionlinesformat[4]{%
    \ifstr{#1}{subsubsection}{%
        \makebox[0pt][r]{#4}%
    }{%
        \makebox[0pt][r]{#3}#4%
    }%
}%

% header & footer
\usepackage{scrlayer-scrpage}
%\lofoot{\today}
%\refoot{\today}
\pagestyle{scrheadings}
%\let\raggedsection\centering

\title{\includegraphics[width = 80mm]{images/logo_crop.png}}
\subtitle{\huge Änderungsdokument}
\author{Lukas Dippon
        \and Jens Kienle
        \and Matthias Noll
        \and Fabian Röpke
        \and Tim Schmidt
        \and Simon Vögele}

\begin{document}

\maketitle
\thispagestyle{empty} % removed page number from title

\pagebreak
\tableofcontents

%%%%%%%%%%%%%%%%%%%
\pagebreak
\section{Entwurf}



%%%%%%%%%%%%%%%%%%%
\pagebreak
\section{Server}



%%%%%%%%%%%%%%%%%%%
\pagebreak
\section{Klient}

\subsection{TreffPunkt}
Für einzelne Daten von primitiven Datentypen bietet es sich an, diese in 
Shared-Preferences zu speichern. Allerdings wird ein Context benötigt, um auf 
diese zuzugreifen. Um Klassen, welche keinen eigenen Context besitzen, trotzdem 
den Zugriff auf Shared-Preferences zu ermöglichen, wurde die TreffPunkt Klasse 
als Superklasse für Application eingeführt. Sie bietet statisch einen Context, 
welcher nicht von irgendwelchen Lebenszyklen abhängig ist.

\subsection{Befehle}
Die interne Umsetzung der Netzwerkbefehle wurde um das Commands Package 
erweitert. Nach dem Command Design Pattern werden nun die einzelnen Server 
Befehle durch Objekte repräsentiert. Die Hauptfunktion der neuen Klassen 
besteht zum einen darin, Objekte für die JSON Serialisierung zu verwalten. Zum 
anderen werden von dort aus aber auch die lokalen Datenbankzugriffe (und 
andere Reaktionen auf Antworten des Servers) gesteuert. So wird vermieden, dass 
die ViewModels auf Reaktionen des Servers warten müssen und die UI einfriert.
Um eine zu enge Kopplung zu verhindern, wird die Ausführung eines Befehls 
weiterhin die jeweilige Methode des RequestEncoders veranlasst\\
Zudem gibt es ein weiteres Package, welches die im Protokoll beschriebenen 
Entitäten enthält. Dieses ist nötig, um eine gewisse Konsistenz zwischen Server 
und Klient zu schaffen.\\
Für die bereits angesprochene JSON Serialisierung wird sowohl Server- , als 
auch Klientseitig die Jackson Library verwendet.

\subsection{Karte}
osmdroid bonuspack hinzugefügt, um clustering zu verwirklichen
marker package im map package: UserMarker, EventMarker, +...

\subsection{Datenbank}
Event entity hat jetzt Gruppen ID
TypeConverter für Integer Sets hinzugefügt
Position entity durch android.location.Location ersetzt

\subsection{Ortung}
GPSProvider nur Service, Map ist selbst auch LocationListener
GPSProviderManager hinzugefügt, um den GPSProvider zu starten

\subsection{Gruppeneinstellungen}
GroupSettingsActivity statt Fragment (BackStack einfacher), ViewModel dazu entfernt/kein eigenes mehr
MemberListAdapter hinzugefügt

\subsection{Über uns}
about package/activity hinzugefügt

\subsection{Login}
LoginFragment hinzufügt

\subsection{Dependency Injection}
AppModule Klasse entfernt (kein Dagger mehr für DI)
Repos keine Singletons, werden einmal in der ViewModelFactory erstellt

\subsection{ViewModels}
util package: ViewCall und State enum
Aufrufe im ViewModel, welche einen Context benötigen, werden in einem SingleLiveEvent State Objekt gespeichert, welches vom Context (Fragment/Activity) observed wird. Dieses kann dann entsprechend auf State-Änderungen reagieren.
ViewModels (nicht die Fragments) erweitern ItemClickListener um alle Nutzerinteraktionen im ViewModel zu behandeln
PagedListAdapter statt normalen Adaptern um unnötiges Laden von ListItems zu vermeiden

\subsection{NavigationDrawer}
TabbedNavigationActivity entfernt, nur noch NavigationActivity
NavigationDrawerFragment nur noch xml, kein Fragment


%%%%%%%%%%%%%%%%%%%
\pagebreak
\section{Protokoll}



\end{document}
